\documentclass[12pt]{article}

\usepackage[margin=1in]{geometry}
\IfFileExists{newtxtext.sty}{
    \usepackage{newtxtext,newtxmath}
}{
    \usepackage{mathptmx}
}
\usepackage{enumitem}
\usepackage{cite}
\usepackage[hidelinks]{hyperref}

\setlist[itemize]{leftmargin=0.28in,itemsep=3pt,topsep=3pt}
\setlist[itemize,2]{leftmargin=0.28in,itemsep=2pt}
\setlength{\parindent}{0pt}
\setlength{\emergencystretch}{2em}

\begin{document}

\begin{center}
    {\large\bfseries  Outline}\\[4pt]
    {\bfseries Testing a One-Dimensional Potential-Barrier Model Using MOSFET Gate-Leakage Data}
\end{center}

\textbf{Research question:} To what extent is the experimentally measured dependence of gate leakage current density on oxide thickness consistent with the prediction of a one-dimensional rectangular potential-barrier model in MOSFETs?

\section*{1. Introduction}
\begin{itemize}
    \item Introduce quantum tunneling as the nonzero probability that a particle with \(E<V_0\) passes through a finite potential barrier \cite{townsend2010}.
    \item Explain why tunneling becomes important when the barrier is extremely thin: the wavefunction has a shorter physical distance over which to decay.
    \item Introduce the MOSFET as a real application:
    \begin{itemize}
        \item The gate controls the conductivity of the silicon channel through an electric field across the gate oxide.
        \item The oxide is an insulating layer, but electrons may tunnel through it when it becomes sufficiently thin.
        \item Gate leakage contributes to power consumption, heating, and reliability problems in scaled devices \cite{green2001review}.
    \end{itemize}
    \item State the research question and explain that the project will test whether a simple quantum model captures the measured \emph{thickness dependence}.
\end{itemize}

\section*{2. Theoretical Background: One-Dimensional Rectangular Barrier}
\begin{itemize}
    \item Define the model:
    \begin{itemize}
        \item Barrier width \(a\), barrier height \(V_0\), and incident electron energy \(E<V_0\).
        \item Three regions: before, inside, and after the barrier.
    \end{itemize}
    \item Set up the time-independent Schr\"odinger equation in each region \cite{townsend2010}.
    \item Describe the expected wavefunctions:
    \begin{itemize}
        \item Incident and reflected traveling waves before the barrier.
        \item Exponentially decaying and growing terms inside the barrier.
        \item A transmitted traveling wave after the barrier.
    \end{itemize}
    \item Apply continuity of the wavefunction and its first derivative at both boundaries.
    \item Derive or present the exact transmission probability \(T\), followed by the high/wide-barrier approximation.
    \item Identify the model's main prediction:
    \begin{itemize}
        \item \(T\) decreases approximately exponentially as \(a\) increases.
    \end{itemize}
    \item State the assumptions: one-dimensional motion, a rectangular barrier, constant \(E\), constant \(V_0\), and no defects or scattering within the oxide.
\end{itemize}

\section*{3. Connecting the Model to a MOSFET}
\begin{itemize}
    \item Map the theoretical quantities to the physical device:
    \begin{itemize}
        \item Barrier width \(a \longleftrightarrow\) oxide thickness \(t_{\mathrm{ox}}\).
        \item Barrier-energy difference \(V_0-E \longleftrightarrow\) an effective energy barrier for an electron entering the oxide.
        \item Transmission probability \(T \longleftrightarrow\) one factor contributing to gate leakage.
    \end{itemize}
    \item Explain the difference between the compared dependent variables:
    \begin{itemize}
        \item \(T\) is a dimensionless probability for an electron.
        \item \(J_G\) is a measured current per unit area and also depends on the number and energy distribution of electrons supplied at the interface.
        \item The investigation will therefore compare slopes, linearity, and normalized trends of \(T\) and \(J_G\).
    \end{itemize}
    \item Explain why the rectangular barrier is only a first approximation:
    \begin{itemize}
        \item An applied electric field tilts the oxide barrier, making it closer to a trapezoidal barrier \cite{simmons1963,green2001review}.
        \item At fixed \(V_G\), changing \(t_{\mathrm{ox}}\) also changes the approximate oxide electric field, so thickness is not the only physical factor changing.
        \item Published MOS tunneling models also account for band bending, electron supply, barrier height, and effective mass \cite{depas1995}.
    \end{itemize}
\end{itemize}

\section*{4. Experimental Data Source}
\begin{itemize}
    \item Use the published NMOSFET data from Green et al. \cite{green1999limits}.
    \item Independent variable: oxide thickness \(t_{\mathrm{ox}}\).
    \item Experimental dependent variable: gate leakage current density \(J_G\) .
    \item Extract the plotted coordinates using a graph digitizer:
    \item Record limitations of the secondary dataset: finite graph resolution, uncertainty in thickness measurements, possible device-to-device variation, and incomplete access to the authors' raw numerical data.
\end{itemize}

\section*{5. Theoretical Calculation}
\begin{itemize}
    \item Use MATLAB to calculate \(T\) over the same oxide-thickness range as the experimental data.
    \item Select and justify values for the effective electron mass and effective barrier energy using semiconductor-tunneling literature \cite{depas1995,green2001review}.
    \item Check units carefully by converting all energies to joules and all thicknesses to meters inside the calculation.
    \item Test the sensitivity of the theoretical slope to plausible literature values of the barrier energy and effective mass using other sets of data.
\end{itemize}

\section*{6. Data Processing and Planned Comparison}
\begin{itemize}
    \item Produce separate raw-scale plots of:
    \begin{itemize}
        \item Theoretical \(T\) against \(t_{\mathrm{ox}}\).
        \item Experimental \(J_G\) against \(t_{\mathrm{ox}}\).
    \end{itemize}
    \item Produce the main comparison plots:
    \begin{itemize}
        \item \(\ln T\) against \(t_{\mathrm{ox}}\).
        \item \(\ln J_G\) against \(t_{\mathrm{ox}}\).
    \end{itemize}
    \item Apply linear regression to both logarithmic relationships and record:
    \begin{itemize}
        \item Slope in \(\mathrm{nm}^{-1}\).
        \item Intercept.
        \item \(R^2\) as a measure of linearity.
        \item Residual pattern.
    \end{itemize}
    \item Normalize both quantities to a common reference thickness if a direct visual comparison of curve shape is needed.
\end{itemize}

\section*{7. Analysis}
\begin{itemize}
    \item Determine whether the model reproduces the experimental exponential thickness dependence.
    \item Interpret the level of agreement:
    \begin{itemize}
        \item Similar logarithmic slopes would support the rectangular-barrier approximation over the selected range.
        \item Similar direction but different slopes would mean that the model is qualitatively useful but quantitatively incomplete.
        \item Curvature or systematic residuals would indicate that a constant rectangular barrier does not fully describe the device.
    \end{itemize}
    \item Discuss possible physical explanations for deviation
    \item Separate evidence of model limitations from uncertainty introduced by graph digitization and the published measurements.
    \end{itemize}

\section*{8. Conclusion}
\begin{itemize}
    \item Answer the research question using both qualitative and quantitative evidence.
    \item State the thickness range over which the model is supported.
    \item Explain what the model captures and which real-device effects it omits.
\end{itemize}

\section*{9. Evaluation and Improvements}
\begin{itemize}
    \item Strengths:
    \begin{itemize}
        \item A simple model gives a clear, testable prediction.
        \item The comparison uses published device measurements with microscopy-based oxide thicknesses.
    \end{itemize}
    \item Limitations and corresponding improvements:
    \begin{itemize}
        \item Rectangular barrier assumption \(\rightarrow\) extend the model to a trapezoidal barrier using a WKB approximation.
        \item Comparison of \(T\) with \(J_G\) \(\rightarrow\) use a fuller current-density model that includes electron supply.
        \item One fixed gate voltage \(\rightarrow\) compare several gate voltages to separate thickness and electric-field effects.
    \end{itemize}
\end{itemize}


\section*{10. Appendices}
\begin{itemize}
    \item MATLAB code.
    \item Full digitized dataset and calculation table.
\end{itemize}

\newpage
\IfFileExists{IEEEtran.bst}{\bibliographystyle{IEEEtran}}{\bibliographystyle{unsrt}}
\bibliography{references}

\end{document}
